% ============================================================================
\chapter{TSS 与 Ring 3：进入用户态}
\label{ch:tss-ring3}
% Stage 10 — 待实现
% ============================================================================

\section{本章目标}

\begin{itemize}
  \item 理解 x86 特权级（Ring 0 -- Ring 3）与特权级切换机制
  \item 实现 Task State Segment（TSS），配置内核栈指针
  \item 在 GDT 中添加 TSS 描述符和用户态代码/数据段
  \item 用 \texttt{iret} 从 Ring 0 跳转到 Ring 3 执行用户代码
\end{itemize}

\section{背景知识：x86 特权级}

% TODO:
% - Ring 0 (内核) / Ring 3 (用户)，Ring 1/2 通常不用
% - CPL (Current Privilege Level): CS 低 2 位
% - DPL (Descriptor Privilege Level): 段/门描述符中的权限
% - RPL (Requested Privilege Level): 段选择子低 2 位
% - 特权级检查规则：CPL <= DPL 才能访问

\subsection{TSS 的作用}

% TODO:
% - Ring 3 → Ring 0 切换时，CPU 从 TSS 读取内核栈指针 (SS0:ESP0)
% - 为什么需要独立的内核栈（安全隔离）
% - TSS 描述符在 GDT 中的格式

\subsection{从 Ring 0 到 Ring 3：\texttt{iret}}

% TODO:
% - 手动在栈上构造 iret 帧: SS3 / ESP3 / EFLAGS / CS3 / EIP3
% - iret 弹出这些值 → CPU 切换到 Ring 3
% - EFLAGS 中 IF=1 保持中断开启

\section{GDT 扩展}

% TODO: 新增 User Code (0x18) / User Data (0x20) / TSS (0x28) 段

\section{TSS 初始化}

% TODO: tss_t 结构体、tss_init()、ltr 指令

\section{跳转到用户态}

% TODO: 汇编辅助函数，构造 iret 帧并执行

\begin{cpustate}
\texttt{iret} 执行后：
\begin{itemize}
  \item CPL 从 0 变为 3（用户态）
  \item CS 指向用户代码段，SS 指向用户数据段
  \item ESP 切换到用户态栈
  \item 内核页表中 PTE\_USER 位控制哪些页用户态可访问
\end{itemize}
\end{cpustate}

\section{验证}

% TODO: 用户态执行一条指令 → 触发 syscall 或 GP fault → 回到内核

\section{本章小结}

TSS + iret 让我们能在 Ring 0 和 Ring 3 之间切换。
下一章加载 ELF 文件，让用户态执行真正的程序。
